% ભાગ: ગણો-વિધેયો-સંબંધો
% પ્રકરણ: અંકગણિતીકરણ
% વિભાગ: કોશી

\documentclass[../../../include/open-logic-section]{subfiles}

\begin{document}

\olfileid{sfr}{arith}{cauchy}

\olsection{પરિશિષ્ટ: કોશી શ્રેણીઓ તરીકે વાસ્તવિક સંખ્યાઓ}

\olref[cuts]{sec}માં આપણે વાસ્તવિક સંખ્યાઓની રચના ડેડેકિન્ડ કાપો
તરીકે કરી. આ વિભાગમાં એક વૈકલ્પિક રચના સમજાવીશું. તે કોશીએ
(જેને હવે આપણે કોશી શ્રેણી કહીએ છીએ તેની) આપેલી વ્યાખ્યા પર આધારિત
છે; પરંતુ આ વ્યાખ્યાનો ઉપયોગ કરીને વાસ્તવિક સંખ્યાઓની \emph{રચના}
કરવાનું શ્રેય ઓગણીસમી સદીના બીજા લેખકોને, ખાસ કરીને વાયરસ્ટ્રાસ,
હાઇને, મેરે અને કૅન્ટૉરને જાય છે. (સરસ ઇતિહાસ માટે
\citeauthor{OConnorRobertson:RN} \citeyear{OConnorRobertson:RN} જુઓ.)

ઓગણીસમી સદી સુધી પહોંચતાં પહેલાં સિમોન સ્ટેવિન (1548--1620)નો વિચાર
કરવો યોગ્ય છે. ટૂંકમાં, સ્ટેવિને સમજ્યું કે દરેક વાસ્તવિક સંખ્યાને
તેના દશાંશ વિસ્તરણ દ્વારા વિચારી શકાય. તેથી $\sqrt{2}$ જેવી અસંમેય
સંખ્યાનું પણ સરસ દશાંશ વિસ્તરણ છે, જે આ રીતે શરૂ થાય છે:
\[
	1.41421356237\ldots
\]
ગણસિદ્ધાંતમાં દશાંશ વિસ્તરણોના પ્રતિરૂપ બનાવવા બહુ સરળ છે: તેમને
$d \colon \Nat \to \Nat$ વિધેયો તરીકે વિચારો, જ્યાં $d(n)$ આપણને
રસ હોય તે $n$મું દશાંશ સ્થાન છે. પછી દશાંશ ચિહ્નની પહેલાં આવતા
વાસ્તવિક સંખ્યાના ભાગને (અહીં માત્ર $1$ને) સંભાળવા વ્યાખ્યામાં થોડો
ફેરફાર કરવો પડશે. દાખલા તરીકે $0.999\ldots = 1$ની ખાતરી કરવા માટે
બીજો ફેરફાર (એક સામ્ય સંબંધ) પણ કરવો પડશે. પરંતુ સ્ટેવિનની રીત મુજબ
ગણસિદ્ધાંતમાં વાસ્તવિક સંખ્યાઓની પૂરેપૂરી ચુસ્ત રચના આપવી મુશ્કેલ નથી.

અહીં સ્ટેવિન આપણો મુખ્ય વિષય નથી. (સ્ટેવિન વિશે વધુ માટે
\citealt{KatzKatz2012} જુઓ.) પરંતુ આ રહ્યો તેની નજીકનો એક વિચાર.
$\sqrt{2}$ના દશાંશ વિસ્તરણને સીધું લેવાને બદલે વધુ ને વધુ ચોક્કસ
વિસ્તરણો લઈને $\sqrt{2}$નાં વધુ ને વધુ ચોક્કસ સંમેય આસન્ન મૂલ્યોની
એક \emph{શ્રેણી} વિચારી શકાય:
\[
1, 1.4, 1.414, 1.4142, 1.41421,\ldots
\]
વાસ્તવિક સંખ્યાઓને ``વધુ ને વધુ સારા આસન્ન મૂલ્યો'' દ્વારા વિચારી
શકાય એ વિચાર આપણને અંતર્દૃષ્ટિઓના બીજા ક્રમનો આધાર આપે છે (જેમ
$\Int$માંથી $\Rat$ અથવા $\Nat$માંથી $\Int$ રચતી વખતે મળેલી
અંતર્દૃષ્ટિઓ). મૂળભૂત અંતર્દૃષ્ટિઓ આ છે:
\begin{enumerate}
	\item દરેક વાસ્તવિક સંખ્યાને (કદાચ અનંત) દશાંશ વિસ્તરણ તરીકે લખી
	શકાય છે.
	\item (કદાચ અનંત) દશાંશ વિસ્તરણમાં સમાયેલી માહિતીને સંમેય સંખ્યાઓની
	શ્રેણીથી પણ એટલી જ રીતે રજૂ કરી શકાય છે.
	\item સંમેય સંખ્યાઓની શ્રેણીને $\Nat$માંથી $\Rat$માં જતા વિધેય
	તરીકે વિચારી શકાય; શ્રેણીની $n$મી સંમેય સંખ્યાને $f(n)$ લો.
\end{enumerate}
અલબત્ત, $\Nat$માંથી $\Rat$માં જતું \emph{કોઈ પણ} વિધેય આપણને વાસ્તવિક
સંખ્યા આપશે એવું નથી. દાખલા તરીકે, આ વિધેય વિચારો:
\[
	f(n) =\begin{cases}
			1 & \text{જો }n\text{ અયુગ્મ હોય}\\
			0 &\text{જો }n\text{ યુગ્મ હોય}
		\end{cases}
\]
અહીં મૂળ ચિંતા એ છે કે $0,1,0,1,0,1,0,\ldots$ શ્રેણી કોઈ વાસ્તવિક
સંખ્યા તરફ ``નજીક જતી'' જણાતી નથી. તેથી કોઈ વાસ્તવિક સંખ્યા તરફ
નજીક જતી શ્રેણીઓ જ વિચારવામાં આવે એની ખાતરી કરવા માટે આપણે કોઈ
\emph{લક્ષ} ધરાવતી શ્રેણીઓ પૂરતું ધ્યાન મર્યાદિત કરવું પડશે.

લક્ષનો વિચાર આપણે \olref[his][set][limits]{sec}માં જોઈ ચૂક્યા છીએ.
પરંતુ ત્યાં લીધેલી વ્યાખ્યા \emph{બરાબર} એ જ રીતે અહીં લઈ શકીએ નહીં.
ત્યાં ``$(\forall \epsilon>0)$'' અભિવ્યક્તિમાં વાસ્તવિક સંખ્યાઓ પરનું
પરિમાણીકરણ ગર્ભિત હતું; અને આપણે વાસ્તવિક સંખ્યાઓ પરનાં વિધેયોના
લક્ષ વિચારતા હતા. એટલે એ વ્યાખ્યાનો ઉપયોગ કરીએ તો જે વાસ્તવિક સંખ્યાઓની
\emph{રચના} કરવાનું આપણું લક્ષ છે તેને પહેલાંથી જ માની લઈએ. સદભાગ્યે,
લક્ષના અત્યંત નજીકના એક વિચારથી કામ ચાલી શકે છે.
\begin{defn}\ollabel{def:CauchySequence}
  વિધેય $f: \Nat \to \Rat$ \emph{કોશી શ્રેણી} હોય ત્યારે અને માત્ર
  ત્યારે જ જ્યારે દરેક ધન $\epsilon \in \Rat$ માટે $(\exists \ell \in
  \Nat)(\forall m, n > \ell)|f(m) - f(n)| < \epsilon$.
\end{defn}

લક્ષનો સામાન્ય વિચાર પહેલાં જેવો જ છે: તમને ($\epsilon$થી મપાતું)
અમુક ચોકસાઈનું સ્તર જોઈતું હોય તો જોવા માટે એક ``વિસ્તાર'' (એટલે કે
$\ell$ કરતાં મોટું કોઈ પણ આગત) મળે છે. આપણી શ્રેણી $1$, $1.4$,
$1.414$, $1.4142$, $1.41421$\ldots{}ને લક્ષ છે એવું જોવું પણ સરળ છે:
$\sqrt{2}$નું $\nicefrac{1}{10^{n}}$ જેટલી ત્રુટિની અંદર આસન્ન મૂલ્ય
જોઈતું હોય તો $n$મા ઘટક પછીના કોઈ પણ ઘટકને જુઓ.

ત્યારે પહેલો સ્વાભાવિક વિચાર એવો આવે કે કોઈ પણ કોશી શ્રેણી બસ એક
વાસ્તવિક સંખ્યા \emph{છે}. પરંતુ $\Int$ અને $\Rat$ની રચનાઓની જેમ આ
વિચાર બહુ ભોળો થશે: આપેલી કોઈ એક વાસ્તવિક સંખ્યાને ઘણી જુદી કોશી
શ્રેણીઓ સૂચવે છે. આ જોવાની એક સરળ રીત આ છે. કોશી શ્રેણી~$f$ આપેલી
હોય ત્યારે $g$ને $f$ જેવું જ વિધેય વ્યાખ્યાયિત કરો, માત્ર $g(0)\neq
f(0)$ રાખો. બંને શ્રેણીઓ પહેલા ઘટક પછી સર્વત્ર એકમત હોવાથી આપણે
(અંતે) એમ કહેવા ઇચ્છીશું કે \olref{def:CauchySequence}માં વપરાયેલા
અર્થમાં તેમને એક જ લક્ષ છે અને તેથી તેઓ એક જ વાસ્તવિક સંખ્યાને
``વ્યાખ્યાયિત'' કરે છે. તેથી આ કોશી શ્રેણીઓને ખરેખર એક જ વાસ્તવિક
સંખ્યા તરીકે વિચારવી જોઈએ.

આથી ફરી કોશી શ્રેણીઓ પર સામ્ય સંબંધ વ્યાખ્યાયિત કરીને વાસ્તવિક
સંખ્યાઓને સામ્ય \emph{વર્ગો} સાથે ઓળખવી પડશે. પહેલાં લક્ષમાં $0$
તરફ અભિસરતા વિધેયનો વિચાર જોઈએ. કોઈ પણ $h : \Nat \to \Rat$ માટે
કહો કે \emph{$h$ લક્ષમાં $0$ તરફ અભિસરે છે} ત્યારે અને માત્ર ત્યારે
જ્યારે દરેક ધન $\epsilon \in \Rat$ માટે $(\exists \ell \in
\Nat)(\forall n > \ell)|h(n)| < \epsilon$.\footnote{$\lim_{x
\mathord{\rightarrow}\infty}f(x) = 0$ની વ્યાખ્યા સાથે સરખાવો;
\olref[his][set][limits]{sec} જુઓ.} વધુમાં, $f$ અને $g$ એ
$\Nat \to \Rat$ વિધેયો હોય ત્યારે $(f-g)(n) = f(n) - g(n)$ લો. હવે
વ્યાખ્યા આપો:
\[
	f \Realequiv g \text{ ત્યારે અને માત્ર ત્યારે જ જ્યારે $(f-g)$ લક્ષમાં $0$ તરફ અભિસરે છે}.
\]
$\Realequiv$ સામ્ય સંબંધ છે એ ચકાસવું પડે; અને તે છે. પછી ઇચ્છીએ તો
વાસ્તવિક સંખ્યાઓને $\Nat \to \Rat$ જતી બધી કોશી શ્રેણીઓના
$\Realequiv$ તળેના સામ્ય વર્ગો તરીકે વ્યાખ્યાયિત કરી શકીએ.
\sourcecorrection{OLARI-003}{મૂળની \texttt{cauchy.tex}, પંક્તિઓ
84--93માં વાસ્તવિક સંખ્યાઓને સામ્ય \emph{સંબંધો} સાથે ઓળખવાનું કહેવાયું
છે, જોકે પછીનું વાક્ય યોગ્ય રીતે સામ્ય વર્ગો વાપરે છે. પંક્તિઓ
152--161નાં પ્રમેય અને પ્રશ્નમાં પણ શ્રેણીઓ કહેવાને બદલે એમના સામ્ય
વર્ગો અભિપ્રેત છે. અહીં ત્રણેય સ્થળે એ વર્ગો સ્પષ્ટ કર્યા છે.}

\begin{prob}
દરેક $n$ માટે $f(n) = 0$ લો. $g(n) = \frac{1}{(n+1)^2}$ લો. બંને
કોશી શ્રેણીઓ છે અને ખરેખર બંને વિધેયોનું લક્ષ $0$ છે, જેથી
$f \sim_\Real g$ પણ થાય છે, તે બતાવો.
\end{prob}

આ કરી લીધા પછી હંમેશની જેમ $f$ને !!a{element} તરીકે ધરાવતા સામ્ય
વર્ગ માટે $\equivrep{f}{\Realequiv}$ લખીશું. પરંતુ વાંચવામાં સરળતા
રહે તે માટે હવે પછી નીચેનો સૂચક છોડીને માત્ર $\equivrep{f}{}$ લખીશું.
દરેક $q \in \Rat$ માટે $q_{\Real} = \equivrep{c_{q}}{}$ પણ ઠરાવીએ
છીએ, જ્યાં દરેક $n \in \Nat$ માટે $c_q(n) = q$ એવું અચળ વિધેય $c_{q}$
છે. પછી વાસ્તવિક સંખ્યાઓ પરના મૂળભૂત સંબંધો અને ક્રિયાઓ આ રીતે
વ્યાખ્યાયિત કરીએ છીએ:
\begin{align*}
	\equivrep{f}{} + \equivrep{g}{} &= 	\equivrep{(f + g)}{} \\
	\equivrep{f}{} \times 	\equivrep{g}{} &= \equivrep{(f \times g)}{}
\end{align*}
જ્યાં $(f + g)(n) = f(n) + g(n)$ અને $(f \times g)(n) = f(n) \times
g(n)$. અલબત્ત, $f$ અને $g$ કોશી શ્રેણીઓ હોય ત્યારે $(f + g)$,
$(f-g)$ અને $(f\times g)$માંનું દરેક કોશી શ્રેણી છે એ પણ ચકાસવું પડે;
તેઓ છે, અને એ કામ તમારા માટે છોડી દઈએ છીએ.

છેલ્લે, ક્રમનો ખ્યાલ વ્યાખ્યાયિત કરીએ. $\equivrep{f}{}$ને
\emph{ધન} કહો ત્યારે અને માત્ર ત્યારે જ જ્યારે $\equivrep{f}{}\neq
0_\Real$ અને $(\exists \ell \in \Nat)(\forall n > \ell)0 < f(n)$
બંને થાય. પછી $\equivrep{f}{} < \equivrep{g}{}$ ત્યારે અને માત્ર
ત્યારે જ જ્યારે $\equivrep{(g - f)}{}$ ધન હોય. આ સુવ્યાખ્યાયિત છે
(એટલે કે સામ્ય વર્ગમાંથી પસંદ કરેલા ``પ્રતિનિધિ'' વિધેય પર આધાર
રાખતું નથી) એ ચકાસવું પડશે.
\sourcecorrection{OLARI-004}{મૂળની \texttt{cauchy.tex}, પંક્તિઓ
137--139માં વાસ્તવિક સંખ્યાના સામ્ય વર્ગની ધનતા વ્યાખ્યાયિત કરતી
વખતે તેને સંમેય શૂન્ય \(0_\Rat\)થી જુદો કહ્યો છે. અહીં યોગ્ય પ્રકારનું
વાસ્તવિક શૂન્ય \(0_\Real\) વાપર્યું છે.}
%\begin{align*}
%	\equivrep{f}{} \leq \equivrep{g}{} \text{ ત્યારે અને માત્ર ત્યારે જ જ્યારે }&\equivrep{f}{} = \equivrep{g}{} \text{ અથવા }\\
%	&(\exists \ell \in \Nat)(\forall n \in \Nat)(\ell < n \lif f(n) \leq g(n))
%\end{align*}
પરંતુ આ ચકાસ્યા પછી આ વ્યાખ્યાઓ યોગ્ય બીજગાણિતિક ગુણધર્મો આપે છે
એ બતાવવું ઘણું સરળ છે; એટલે કે:
\begin{thm}\ollabel{thm:cauchyorderedfield}
કોશી શ્રેણીઓના સામ્ય વર્ગો ક્રમિત ક્ષેત્ર બનાવે છે.
\end{thm}

\begin{proof}
કસરત.
\end{proof}

\begin{prob}
કોશી શ્રેણીઓના સામ્ય વર્ગો ક્રમિત ક્ષેત્ર બનાવે છે તે સાબિત કરો.
\end{prob}

આ રીતે રચેલી વાસ્તવિક સંખ્યાઓ પૂર્ણતા ગુણધર્મ ધરાવે છે એવું સાબિત
કરવું વધુ મુશ્કેલ છે, તેથી આપણે તેની સાબિતી આપીશું.

\begin{thm}
ઉચ્ચસીમા ધરાવતા કોશી શ્રેણીઓના દરેક અરિક્ત ગણને ન્યૂનતમ ઉચ્ચસીમા છે.
\end{thm}

\begin{proof}[સાબિતીની રૂપરેખા] ઉચ્ચસીમા ધરાવતો કોશી શ્રેણીઓનો કોઈ
અરિક્ત ગણ $S$ લો. તેથી કોઈ $p \in \Rat$ એવો છે કે $p_{\Real}$ એ
$S$ની ઉચ્ચસીમા છે. $r \in S$ લો; ત્યારે કોઈ $q \in \Rat$ એવો છે
કે $q_{\Real} < r$. તેથી $S$ની ન્યૂનતમ ઉચ્ચસીમા અસ્તિત્વ ધરાવતી હોય
તો તે $q_\Real$ અને $p_\Real$ની વચ્ચે (અંતિમ મૂલ્યો સહિત) છે.

નીચેથી અને ઉપરથી એકસાથે નજીક જઈને આપણે ન્યૂનતમ ઉચ્ચસીમા તરફ
સંકુચિત થઈશું. ખાસ કરીને, $f$ ઉપરથી અને $g$ નીચેથી ન્યૂનતમ
ઉચ્ચસીમા તરફ સંકુચિત થાય એ હેતુથી બે વિધેયો $f, g \colon \Nat \to
\Rat$ વ્યાખ્યાયિત કરીએ છીએ. શરૂઆતમાં આ વ્યાખ્યાઓ આપીએ:
\begin{align*}
	f(0) &= p \\
	g(0) &= q
\end{align*}
પછી $a_n = \frac{f(n) + g(n)}{2}$ હોય ત્યારે લો:\footnote{આ
આવર્તક વ્યાખ્યા છે. પરંતુ આવર્તક વ્યાખ્યાઓ વાજબી છે એવું વિચારવાનું
કોઈ કારણ આપણે \emph{હજી} આપ્યું નથી.}
\begin{align*}
	f(n+1) &=
	\begin{cases}
		a_n &\text{જો }(\forall h \in S)\equivrep{h}{} \leq (a_n)_\Real\\
		f(n)&\text{અન્યથા}
	\end{cases}\\
	g(n+1) &=
	\begin{cases}
		a_n &\text{જો }(\exists h \in S)\equivrep{h}{} \geq (a_n)_\Real\\
		g(n) &\text{અન્યથા}
	\end{cases}
\end{align*}
$f$ અને $g$ બંને કોશી શ્રેણીઓ છે. (આ ચકાસવું ઘણું સરળ છે, પણ આપણે
તે કસરત તરીકે છોડીએ છીએ.) $(f-g)$ વિધેય લક્ષમાં $0$ તરફ અભિસરે છે,
કારણ કે દરેક પગલે $f$ અને $g$ વચ્ચેનો તફાવત અડધો થાય છે. તેથી
$\equivrep{f}{} = \equivrep{g}{}$.

પહેલાં બતાવીએ કે $\equivrep{f}{}$ એ $S$ની ઉચ્ચસીમા છે, એટલે કે
$(\forall h \in S)\equivrep{h}{} \leq \equivrep{f}{}$.
(આગળ વધતાં આપણે \olref{thm:cauchyorderedfield}નો ઉપયોગ કરીશું.)
$h \in S$ લો અને પ્રતિષેધ માટે ધારો કે $\equivrep{f}{} <
\equivrep{h}{}$, જેથી $0_\Real < \equivrep{(h-f)}{}$. $f$ એકધારી રીતે
ઘટતી કોશી શ્રેણી હોવાથી કોઈ $n \in \Nat$ એવો છે કે
$\equivrep{(c_{f(n)} - f)}{} < \equivrep{(h-f)}{}$. તેથી:
\[
	(f(n))_\Real = \equivrep{c_{f(n)}}{} < \equivrep{f}{} + \equivrep{(h-f)}{} = \equivrep{h}{},
\]
જે રચના મુજબ $\equivrep{h}{} \leq (f(n))_\Real$ હોવાનો વિરોધાભાસ છે.

હવે બતાવીએ કે $\equivrep{f}{} = \equivrep{g}{}$ એ $S$ની
\emph{ન્યૂનતમ} ઉચ્ચસીમા છે. કોઈ કોશી શ્રેણી $j$ લો અને ધારો કે
$\equivrep{j}{} < \equivrep{g}{}$. ઉપરની જેમ વિચારતાં ($g$
\emph{વધતી} છે એ હકીકતનો ઉપયોગ કરીને), કોઈ $n \in \Nat$ એવો છે કે
$\equivrep{j}{} < (g(n))_\Real$. પરંતુ રચના મુજબ કોઈ $h \in S$ એવો
છે કે $(g(n))_\Real \leq \equivrep{h}{}$, તેથી $\equivrep{j}{} <
\equivrep{h}{}$ અને એટલે $\equivrep{j}{}$ એ $S$ની ઉચ્ચસીમા નથી.
\end{proof}

\end{document}
